\documentclass[11pt]{article}

%------------------------------------------------
% Packages
%------------------------------------------------
\usepackage[a4paper,margin=2.5cm]{geometry}
\usepackage{amsmath,amssymb,amsthm}
\usepackage{mathtools}
\usepackage{enumitem}
\usepackage{hyperref}

\hypersetup{
    colorlinks=true,
    linkcolor=blue,
    citecolor=blue,
    urlcolor=blue
}

%------------------------------------------------
% Title Information
%------------------------------------------------

\title{\bfseries Derivative-Free Bilevel Learning for Inverse Problems: From Stochastic Gradient Estimation to Multilevel Methods}

\author{
Neil K. Chada\thanks{Department of Mathematics, City University of Hong Kong, Hong Kong SAR, \texttt{neilchada123@gmail.com}}
\and
Mohammad S. Salehi\thanks{Independent scholar, London, UK, \texttt{sadeghsalehi1997@gmail.com}}
\and 
Mathias Staudigl\thanks{Institute of Mathematics, University of Mannheim, Germany, \texttt{m.staudigl@uni-mannheim.de}}
\and 
\underline{Simon Weissmann}\thanks{Institute of Mathematics, University of Mannheim, Germany, \texttt{simon.weissmann@uni-mannheim.de}}
}


\date{}

%------------------------------------------------

\begin{document}

\maketitle

\vspace{-1em}

\noindent
\textbf{Keywords:}
Bilevel optimization, derivative-free optimization, multilevel Monte Carlo, Gaussian smoothing, inverse problems.

\vspace{0.5em}


\vspace{1em}

%%%%%%%%%%%%%%%%%%%%%%%%%%%%%%%%%%%%%%%%%%%%%%%%%%
\section*{Abstract}

We consider bilevel learning problems arising in inverse problems, where regularization parameters are optimized from data and each upper-level function evaluation requires the approximate solution of a lower-level reconstruction problem. When derivatives of the lower-level solution map are unavailable or undesirable, zeroth-order methods provide a flexible alternative, but can become computationally expensive due to increasingly accurate inner solves. We propose a multilevel Monte Carlo approach for derivative-free bilevel optimization that combines Gaussian smoothing with a hierarchy of lower-level accuracies. The resulting estimator uses many inexpensive coarse solves and only few accurate fine solves. We derive convergence and complexity guarantees and identify the properties of the lower-level solver that determine whether multilevel sampling yields a genuine computational speedup. Numerical experiments for imaging inverse problems illustrate the potential gains. This

%%%%%%%%%%%%%%%%%%%%%%%%%%%%%%%%%%%%%%%%%%%%%%%%%%
%\vspace{1em}
%
%\noindent
%\textbf{Acknowledgements.}
%This work was supported by ...
%
%%%%%%%%%%%%%%%%%%%%%%%%%%%%%%%%%%%%%%%%%%%%%%%%%%%
%\vspace{1em}
%
%\begin{thebibliography}{9}
%
%\bibitem{Attouch2013}
%H.~Attouch, J.~Bolte and B.~Svaiter.
%\newblock Convergence of descent methods for semi-algebraic and tame problems.
%\newblock {\em Mathematical Programming}, 137 (2013), 91--129.
%
%\bibitem{Nesterov2018}
%Y.~Nesterov.
%\newblock {\em Lectures on Convex Optimization}.
%\newblock Springer, 2018.
%
%\end{thebibliography}

\end{document}